\subsection{The prime curve does not satisfy the cubic bound}

Let
\[
 S_n^{(1)}=-\pi(X_n)
\]
be the contribution of the curve $c=1$, and put
\[
 S_n^{(\ge2)}=S_n-S_n^{(1)}.
\]
Then
\begin{equation}
 \|S\|^2
 =
 \|S^{(1)}\|^2
 +
 \|S^{(\ge2)}\|^2
 +
 2\langle S^{(1)},S^{(\ge2)}\rangle.
 \label{eq:prime-rest-energy}
\end{equation}

\begin{table}[htbp]
\centering
\caption{Prime-curve versus remaining-curve energy decomposition.}
\label{tab:prime-curve-decomposition}
\scriptsize
\setlength{\tabcolsep}{3.2pt}
\begin{tabular}{rcccccc}
\toprule
$N$ &
$V/N^3$ &
$V_{c=1}/N^3$ &
$V_{c\ge2}/N^3$ &
$2\langle c=1,c\ge2\rangle/N^3$ &
$V/V_{c=1}$ &
correlation\\
\midrule
128  & 0.008842 & 51.225450 & 50.965559 & $-102.182168$ & $1.7260\!\times10^{-4}$ & $-0.999916714$\\
256  & 0.011671 & 146.590116 & 146.035583 & $-292.614028$ & $7.9618\!\times10^{-5}$ & $-0.999961911$\\
512  & 0.007442 & 441.694929 & 441.211486 & $-882.898974$ & $1.6848\!\times10^{-5}$ & $-0.999991721$\\
1024 & 0.009853 & 1382.291036 & 1381.948964 & $-2764.230146$ & $7.1282\!\times10^{-6}$ & $-0.999996443$\\
2048 & 0.008497 & 4450.731250 & 4450.668331 & $-8901.391084$ & $1.9091\!\times10^{-6}$ & $-0.999999045$\\
4096 & 0.009974 & 14652.537064 & 14655.994904 & $-29308.521994$ & $6.8072\!\times10^{-7}$ & $-0.999999667$\\
\bottomrule
\end{tabular}
\end{table}

At $N=4096$, the prime curve has energy more than $1.46\times10^4N^3$, while the total residual is approximately $9.97\times10^{-3}N^3$.  The squared norm of the total is only $6.81\times10^{-7}$ of the prime-curve squared norm.  This is direct numerical evidence that the central phenomenon is cancellation \emph{between} cofactor curves, not cubic behavior of the prime curve in isolation.

